% Challenge dossier: VI.01.C2
\subsection*{Challenge VI.01.C2 --- Multivariate interpolation on an arbitrary finite set}
\label{chal:vi01-c2}

\paragraph{Challenge.}
Let \(F\subseteq k^n\) be finite.  Prove that every function
\[
\varphi:F\to k
\]
is the restriction of a polynomial in \(k[x_1,\dots,x_n]\).  Construct the polynomial
explicitly.

\ifdefined\FullProblemDossiers
\paragraph{Why this challenge matters.}
Polynomial functions can separate any finite collection of affine points, over any field.
This is a concrete finite-set version of the algebra--geometry correspondence.

\paragraph{Guided hints.}
For distinct points \(p,q\in F\), choose a coordinate in which they differ and construct a
linear polynomial that is \(1\) at \(p\) and \(0\) at \(q\).  Multiply these factors over
all \(q\neq p\).

\paragraph{Solution.}
Fix \(p\in F\).  For every \(q\in F\setminus\{p\}\), choose an index \(i(q)\) such that
\(p_{i(q)}\neq q_{i(q)}\), and define
\[
\ell_{p,q}(x)=
\frac{x_{i(q)}-q_{i(q)}}{p_{i(q)}-q_{i(q)}}.
\]
Then \(\ell_{p,q}(p)=1\) and \(\ell_{p,q}(q)=0\).  Put
\[
\delta_p(x)=\prod_{q\in F\setminus\{p\}}\ell_{p,q}(x).
\]
Thus
\[
\delta_p(p)=1,
\qquad
\delta_p(q)=0\quad(q\neq p).
\]
For a prescribed function \(\varphi:F\to k\), the polynomial
\[
f(x)=\sum_{p\in F}\varphi(p)\delta_p(x)
\]
satisfies \(f|_F=\varphi\).

\paragraph{Extensions.}
Use the indicator polynomials to prove directly that every finite subset of affine space is
algebraic and to derive a Chinese-remainder description of its coordinate algebra.

\paragraph{Provenance.}
New challenge; a multivariable Lagrange interpolation construction.
\fi
